\documentclass[11pt,a4paper]{ctexart}
\usepackage[teacher]{ipara}
\tcbuselibrary{breakable}
\usepackage{multicol}
\usepackage{pgfplots}
\pgfplotsset{compat=1.18}
\usetikzlibrary{arrows.meta,positioning,decorations.pathreplacing,shapes.geometric,calc}

\pagestyle{fancy}
\fancyhf{}
\lhead{\small 高中数学方法专题：代数条件的几何意义}
\rhead{\small 教师版（恒成立、定点圆轨迹、三角形面积模型）}
\cfoot{\small 第 \thepage\ 页\quad 共 \zpageref{LastPage} 页}
\renewcommand{\headrulewidth}{0.4pt}
\renewcommand{\footrulewidth}{0pt}

\renewcommand{\teacherproblem}[2]{%
  \teachquestionheader{例题 #1\quad #2}
}

\begin{document}

\begin{center}
  {\zihao{2}\bfseries 高中数学专题讲义 \quad 代数条件的几何意义}\\[3mm]
  {\Large \bfseries 从数到形的直观转化与降维打击}\\[2mm]
  {\large 适用：高一/高二拔高 / 一轮复习专题课 / 120 分钟教案}\\[2mm]
  \rule{0.85\linewidth}{0.5pt}
\end{center}

\begin{teachbox}{专题设计定位}
本专题重在引导学生克服“机械代数运算”的思维惰性，建立\textbf{“代数式与几何图形的翻译映射”}。\\
高中数学压轴题中，很多看似繁复的代数不等式、分式最值、多参约束或根式复合函数，如果纯靠代数通法（如导数、判别式、三角代换等）硬算，不仅运算量巨大，且极易出错。然而，若能洞察其背后的几何本质——如抛物线恒非负对应可行域边界、线性分式对应两点连线斜率、垂足投影对应直径圆轨迹、根式勾股结构对应直角三角形边长与面积，就能将代数最值降维为几何临界状态（如相切、共线、垂直），从而实现“一眼看穿答案，极简完成算术”。\\
本教案旨在引导教师在 120 分钟内，通过三类最典型的问题，帮助学生打破数形界限，培养几何直觉与临界思维。
\end{teachbox}

% ==================== 教学大纲 ====================
\section{120分钟教学大纲与规划}

\begin{center}
\renewcommand{\arraystretch}{1.2}
\begin{tabularx}{0.98\linewidth}{>{\bfseries}p{2.2cm} c p{4.2cm} X}
\toprule
核心板块 & 建议时间 & 核心教学内容 & 教学目标与关键动作 \\
\midrule
核心公式与引子 & 25min & 8 大核心公式与题感诊断 & 快速测试学生对代数式几何意义的敏感度，建立直观印象。 \\
题型一：斜率最值 & 25min & 恒成立不等式与双曲线切线 & 讲解如何将二次式恒成立转化为双曲线可行域，分式转化为斜率，通过相切求最值。 \\
题型二：圆轨迹 & 25min & 偶函数零点与垂足圆轨迹 & 讲解利用偶函数零点对称性求参，识别直线族定点，将垂直转化为直径圆轨迹。 \\
题型三：面积模型 & 25min & 根式勾股与三角形正弦面积 & 讲解如何把双根式和自变量转化为直角三角形高度与底边，利用面积最值破局。 \\
方法联系与思考 & 20min & 通法取舍、变式与实战做题启发 & 总结四类几何意义的横向迁移，对比代数法与几何法的边界，提炼实战启发。 \\
\bottomrule
\end{tabularx}
\end{center}

\newpage

% ==================== 第一部分 ====================
\section{核心引子与题感诊断}

\subsection{1. 本节课用到的核心公式、定理或二级结论}

\begin{keybox}{1.1 二次函数恒非负}
设 $q(x) = ax^2 + bx + c$。若对任意 $x \in \mathbb{R}$ 均有 $q(x) \ge 0$，且 $q(x)$ 为真正的二次函数，则必须满足：
$$a > 0,\qquad \Delta = b^2 - 4ac \le 0$$
在本节第一题中，已知 $b < 0$。若 $a = 0$，则 $bx + c$ 不可能在整个实数集上恒非负（一次函数在 $\mathbb{R}$ 上值域为 $\mathbb{R}$），因此本题确有 $a > 0$。于是：
$$ac \ge \frac{b^2}{4} > 0$$
进而可得 $c > 0$。
\end{keybox}

\begin{keybox}{1.2 斜率的坐标形式}
两点 $P(x_1, y_1)$ 与 $Q(x_2, y_2)$ 满足 $x_1 \ne x_2$ 时，直线 $PQ$ 的斜率为：
$$k_{PQ} = \frac{y_2 - y_1}{x_2 - x_1}$$
当目标代数式能够整理为“纵坐标之差除以横坐标之差”的形式时，可尝试将代数最值转化为平面区域中两点连线的斜率最值。
\end{keybox}

\begin{keybox}{1.3 直线与曲线相切}
将直线与曲线联立消元后得到一元二次方程。若直线与曲线相切，则该二次方程有两个相等的实根，因此判别式满足：
$$\Delta = 0$$
但需要特别注意，由 $\Delta = 0$ 求出多个候选切线斜率值后，必须检查实际切点属于曲线的哪个分支，以及是否符合原题参数的取值范围（如正负性、第一象限等限制）。
\end{keybox}

\begin{keybox}{1.4 偶函数唯一零点}
若函数 $f(x)$ 是偶函数，则对定义域内任意 $x$ 均有：
$$f(-x) = f(x)$$
因此，若 $x_0 \ne 0$ 是该函数的一个零点，则由对称性，$-x_0$ 也必是零点。\\
所以，偶函数若有且只有一个零点，该零点只能是：
$$x = 0$$
这也是求参时优先代入的临界条件。
\end{keybox}

\begin{keybox}{1.5 投影点的圆轨迹}
设 $P, A$ 为定点，动直线 $l$ 经过定点 $A$，点 $M$ 是点 $P$ 在直线 $l$ 上的投影（即 $PM \perp l$）。\\
由于 $PM \perp AM$，所以 $\angle PMA = 90^\circ$ 恒成立。\\
根据圆周角定理，点 $M$ 的运动轨迹是以线段 $PA$ 为直径的圆。
\end{keybox}

\begin{keybox}{1.6 点到圆上动点的距离}
设圆心为 $O_0$、半径为 $r$，定点为 $N$，圆上动点为 $M$。\\
若定点 $N$ 在圆外（即 $O_0N > r$），则：
$$MN_{\min} = O_0N - r,\qquad MN_{\max} = O_0N + r$$
等号在 $O_0, M, N$ 三点共线时取得。
\end{keybox}

\begin{keybox}{1.7 根式的勾股意义}
当 $x^2 \le R^2$ 时，显然有：
$$x^2 + \left(\sqrt{R^2 - x^2}\right)^2 = R^2$$
因此，自变量 $x$ 与根式 $\sqrt{R^2 - x^2}$ 可以看作是一个斜边长度为 $R$ 的直角三角形的两条直角边。
\end{keybox}

\begin{keybox}{1.8 三角形正弦面积公式}
若三角形两邻边长度分别为 $p, q$，它们的夹角为 $\theta$，则该三角形的面积为：
$$S = \frac{1}{2}pq\sin\theta$$
由于对任意的三角形夹角均有 $\sin\theta \le 1$，所以：
$$S \le \frac{1}{2}pq$$
当且仅当 $\theta = 90^\circ$（即两边互相垂直）时取得最大值。
\end{keybox}

\subsection{2. 看到这些条件，应立刻想到什么}

\begin{itemize}
  \item \textbf{看到“对任意 $x \in \mathbb{R}$，二次式恒非负”}：\\
        优先写出开口向上且判别式不大于零：$a > 0, \Delta \le 0$。若进一步观察到 $a, c$ 的乘积形式，则可将 $ac \ge b^2/4$ 翻译为：平面直角坐标系中，点 $(a,c)$ 位于双曲线 $xy = b^2/4$ 的边界及其外侧（上方）。
  \item \textbf{看到线性分式（如 $\frac{b+c}{a-2b}$）}：\\
        应尝试改写为纵坐标之差与横坐标之差之比，即
        $$\frac{c - (-b)}{a - 2b}$$
        此时，它表示动点 $(a,c)$ 与定点 $(2b,-b)$ 连线的斜率 $k$。
  \item \textbf{看到“有且只有一个零点”且函数为偶函数}：\\
        立即想到零点关于原点对称，唯一零点只能是 $x = 0$，从而优先代入并检验 $f(0) = 0$。
  \item \textbf{看到含参数的直线族（如 $m(x-1) + n(y-1) = 0$）}：\\
        首先寻找其恒过的定点。令 $x-1=0, y-1=0$，得 $x=1, y=1$，即所有直线均经过定点 $A(1,1)$。
  \item \textbf{看到“投影”与“垂直”}：\\
        若动点 $M$ 是定点 $P$ 在绕定点 $A$ 旋转的动直线上的投影，应立即联想到直角和圆，即以 $PA$ 为直径的圆轨迹。
  \item \textbf{看到根式复合形式（如 $\sqrt{a-x^2}$ 与 $\sqrt{1-x^2}$）}：\\
        应联想到两个斜边分别为 $\sqrt{a}$ 和 $1$ 的直角三角形，其中一条直角边均为 $x$。若它们与 $x$ 作乘积，应考虑是否可以通过拼接直角三角形，将代数式还原为更大三角形的面积公式（底乘高）。
\end{itemize}

\subsection{3. 本类题的常见切入点}

\begin{itemize}
  \item \textbf{切入点一：把参数看作平面内的点}：\\
        当题目中出现参数间的不等式约束（如 $ac \ge \frac{b^2}{4}$）时，可以定义点 $Q(a,c)$，从而将复杂的参数约束转化为直观的几何区域（第一象限双曲线边界及上方）。
  \item \textbf{切入点二：构造斜率}：\\
        当目标式含有一次式之比时，尝试配凑成 $\frac{y_2-y_1}{x_2-x_1}$ 的结构，将代数最值转化为动点到定点的斜率临界问题。
  \item \textbf{切入点三：先找动直线的定点}：\\
        动直线方程中参数较多时，不要盲目联立求垂足的坐标表达式，而应先找直线恒过的定点，用“定点所对直角”直接锁定圆轨迹。
  \item \textbf{切入点四：把根式还原为图形边长}：\\
        避开直接求导的繁重运算。若根式内部有平方差结构，通过构造直角三角形将代数式翻译为面积，利用夹角直角时面积最大快速求解。
\end{itemize}

\subsection{4. 教学与备考价值判断}

本节课的核心价值，在于训练学生在面对复杂代数式时，能够超越死板的算术，主动去识别斜率、定点、圆轨迹和三角形面积等几何结构。\\
这套方法提供了极其直观的物理与几何背景，尤其适合处理以下问题：
\begin{enumerate}[label=\arabic*., leftmargin=2em, itemsep=0.1em]
  \item 参数约束条件多但结构规则的高考压轴题；
  \item 代数求导运算量极大，或不易判断单调性的最值问题；
  \item 需要严格筛选“等号成立条件”和“临界位置”的多参最值问题；
  \item 动态投影、动点轨迹与最短距离的解析几何综合题。
\end{enumerate}

\newpage

% ==================== 第二部分 ====================
\section{核心思路提炼与板块重构}

\teacherproblem{一}{恒成立不等式中的斜率与切线}

\begin{exbox}{【例题 1】}
已知对任意 $x \in \mathbb{R}$，均有不等式 $ax^2 + bx + c \ge 0$ 成立，其中 $b < 0$。若存在 $t \in \mathbb{R}$，使得 $(1-t)a + (1+2t)b + 3c = 0$ 成立，求 $t$ 的最小值。
\end{exbox}

\par\addvspace{0.4em}{\centering
\begin{tikzpicture}[scale=0.8, >=stealth]
  % 网格
  \draw[color=gridcolor, thin, step=0.5cm] (-5,-1.5) grid (6,5);
  % 坐标轴
  \draw[->, thick] (-5.5,0) -- (6.5,0) node[right] {$x$ (参数 $a$)};
  \draw[->, thick] (0,-1.8) -- (0,5.5) node[above] {$y$ (参数 $c$)};
  \node[below left] at (0,0) {$O$};
  
  % 可行域阴影 xy >= 1, x>0, y>0
  \fill[gray!15] (0.2, 5) -- plot[domain=0.2:5, samples=100] (\x, {1/\x}) -- (5, 5) -- cycle;
  
  % 边界双曲线 xy = 1
  \draw[very thick, blue, domain=0.18:5.2, samples=100] plot (\x, {1/\x});
  \node[blue, right] at (1.5, 1.2) {$ac = \frac{b^2}{4}$};
  
  % 定点 P(2b, -b) => 设 b=-2 => P(-4, 2)
  \coordinate (P) at (-4,2);
  \fill[red] (P) circle (2pt) node[above left] {$P(2b, -b)$};
  
  % 切点 T(2, 0.5)
  \coordinate (T) at (2,0.5);
  \fill[red] (T) circle (2.2pt) node[above right] {$T(a,c)$};
  
  % 切线 y = -0.25x + 1
  \draw[thick, dashed, red, domain=-5.2:5.2] plot (\x, {-0.25*\x + 1});
  
  \node[above] at (3.5, 3.2) {可行域 $ac \ge \frac{b^2}{4}$};
\end{tikzpicture}\par}\addvspace{0.4em}

\begin{paracol}{2}
\teachblock{标准解答与讲解主线}{
  \textbf{第一步：由恒成立不等式确定参数范围与几何区域}\\
  因为对任意 $x \in \mathbb{R}$ 均有 $ax^2+bx+c \ge 0$，且已知 $b < 0$。\\
  我们首先分析二次项系数 $a$ 的范围：
  \begin{itemize}
    \item 若 $a=0$，则不等式退化为一次不等式 $bx+c \ge 0$。由于已知 $b < 0$，一次函数 $y=bx+c$ 的图像是一条斜向下的直线，其值域为 $\mathbb{R}$。当 $x$ 足够大时，函数值必为负数，无法在整个实数集 $\mathbb{R}$ 上恒非负。因此必有 $a \ne 0$。
    \item 若 $a < 0$，抛物线开口向下，其在 $x \to \infty$ 时函数值趋向于负无穷，不可能在 $\mathbb{R}$ 上恒非负。
  \end{itemize}
  综上，二次项系数必须满足开口向上，即 $a > 0$。\\
  在开口向上的前提下，抛物线要恒处于 $x$ 轴上方或与 $x$ 轴接触，二次方程 $ax^2+bx+c=0$ 必须至多有一个实根（即图像与 $x$ 轴至多有一个交点），因此其判别式必须满足：
  $$\Delta = b^2 - 4ac \le 0$$
  将判别式变形：
  $$4ac \ge b^2 \implies ac \ge \frac{b^2}{4}$$
  因为已知 $a > 0$ 且 $b < 0 \implies b^2 > 0$，所以有：
  $$ac \ge \frac{b^2}{4} > 0$$
  由于两数之积为正数，且其中一个数 $a > 0$，故另一个数必有 $c > 0$。\\
  若我们在平面直角坐标系中以 $a$ 为横坐标、$c$ 为纵坐标，令动点 $Q = (a, c)$，则点 $Q$ 位于第一象限。其满足不等式 $xy \ge b^2/4$，表示点 $Q$ 的可行域为双曲线 $xy = b^2/4$ 及其上方的区域。\\[0.3em]
  \textbf{第二步：代数整理，分离参数 $t$}\\
  我们需要从等式中解出 $t$。将等式展开并进行整理：
  $$(1-t)a + (1+2t)b + 3c = 0 \implies a - ta + b + 2tb + 3c = 0$$
  将含 $t$ 的项移到等式右侧，不含 $t$ 的项保留在左侧：
  $$a + b + 3c = ta - 2tb \implies a + b + 3c = t(a - 2b)$$
  因为已知 $a > 0$ 且 $b < 0$，所以有 $-2b > 0$，从而其和满足 $a-2b > 0$ 恒成立。\\
  由于 $a - 2b$ 不为零，我们可以在等式两边同时除以 $a-2b$，解出 $t$ 的表达式：
  $$t = \frac{a+b+3c}{a-2b}$$
  为了方便进行几何转化，我们拆分分子使其显式出现分母的倍数关系：
  $$t = \frac{(a-2b) + 3b + 3c}{a-2b} = \frac{a-2b}{a-2b} + \frac{3(b+c)}{a-2b} = 1 + 3 \cdot \frac{b+c}{a-2b}$$
  \textbf{第三步：将分式解释为几何斜率}\\
  因为平面上任意两点 $A(x_1, y_1)$ 与 $B(x_2, y_2)$ 连线的斜率公式为 $k = (y_2 - y_1)/(x_2 - x_1)$，我们可将上式中的分式改写为减法形式：
  $$\frac{b+c}{a-2b} = \frac{c - (-b)}{a - 2b}$$
  设定点为 $P = (2b, -b)$。由于已知 $b < 0$，所以其横坐标 $2b < 0$，纵坐标 $-b > 0$，定点 $P$ 恒位于第二象限。\\
  设动点为 $Q = (a, c)$，则定点 $P$ 与动点 $Q$ 连线直线的斜率即为：
  $$k_{PQ} = \frac{c - (-b)}{a - 2b}$$
  因此目标参数为：$t = 1 + 3k_{PQ}$。\\
  由于常数 $1$ 和 $3$ 均为定值，且 $3 > 0$，所以 $t$ 的大小随斜率 $k_{PQ}$ 的增大而增大。求 $t$ 的最小值，完全等价于求直线 $PQ$ 斜率的最小值 $k_{\min}$。\\[0.3em]
  \textbf{第四步：利用切线确定斜率的极值}\\
  定点 $P(2b, -b)$ 是第二象限的一个固定点，动点 $Q(a,c)$ 在第一象限双曲线边界及上方区域运动。\\
  从 $P$ 出发向该可行域引射线，当直线由上向下旋转时，斜率 $k$ 逐渐减小。当斜率降低到极限状态时，直线恰好与双曲线边界 $xy = b^2/4$ 在第一象限相切。\\
  设过定点 $P(2b, -b)$ 的切线斜率为 $k$，则切线方程为：
  $$y - (-b) = k(x - 2b) \implies y = kx - (2k+1)b$$
  将切线方程与双曲线方程 $xy = \frac{b^2}{4}$ 联立，消去 $y$ 整理为关于 $x$ 的方程：
  $$x\left[kx - (2k+1)b\right] = \frac{b^2}{4} \implies kx^2 - (2k+1)bx - \frac{b^2}{4} = 0$$
  这是一元二次方程。因为相切时只有一个交点（即方程有且仅有一个重根），所以方程判别式必须为零：
  $$\Delta = \left[-(2k+1)b\right]^2 - 4 \cdot k \cdot \left(-\frac{b^2}{4}\right) = 0$$
  展开并提取公因式 $b^2$：
  $$(2k+1)^2b^2 + kb^2 = 0 \implies b^2\left[(2k+1)^2 + k\right] = 0$$
  因为 $b < 0 \implies b^2 > 0$，我们可在两边约去非零常数 $b^2$，得到：
  $$(2k+1)^2 + k = 0 \implies 4k^2 + 5k + 1 = 0$$
  这是一个关于 $k$ 的二次方程，因式分解为 $(4k+1)(k+1) = 0$，解得两个切线斜率候选值：
  $$k = -1 \quad \text{或} \quad k = -\frac{1}{4}$$
  \textbf{第五步：切点位置的检验与分支筛选}\\
  我们必须验证这两个斜率对应的实际切点是否落在第一象限（即横坐标 $a > 0$）。\\
  联立方程的重根（即切点的横坐标 $a$）公式为 $x_0 = -B/(2A)$。代入本题二次方程的系数 $A=k, B=-(2k+1)b$ 可得：
  $$x_0 = \frac{(2k+1)b}{2k}$$
  - 当 $k = -1$ 时，切点横坐标为：
    $$x_0 = \frac{[2(-1)+1]b}{2(-1)} = \frac{-b}{-2} = \frac{b}{2}$$
    因为已知 $b < 0$，所以此时切点横坐标 $x_0 < 0$，切点位于第三象限的双曲线分支上，不符合 $a>0$ 的实际题意，舍去。
  - 当 $k = -\frac{1}{4}$ 时，切点横坐标为：
    $$x_0 = \frac{[2(-\frac{1}{4})+1]b}{2(-\frac{1}{4})} = \frac{0.5b}{-0.5} = -b$$
    因为已知 $b < 0$，所以此时切点横坐标 $x_0 = -b > 0$。代入双曲线方程可得纵坐标 $c = \frac{b^2}{4a} = -\frac{b}{4} > 0$，切点落在第一象限内，符合题意。
  因此，斜率的最小值为 $k_{\min} = -\frac{1}{4}$。\\
  从而 $t$ 的最小值为：
  $$t_{\min} = 1 + 3\left(-\frac{1}{4}\right) = \frac{1}{4}$$
  \textbf{第六步：检查等号成立条件}\\
  当参数 $a = -b$ 且 $c = -b/4$ 时，显然满足 $a > 0, c > 0$ 且满足 $ac = b^2/4$ 边界条件，此时等号完全可以取得。\\
  所以，常数 $t$ 的最小值为 $\frac{1}{4}$。
}
\switchcolumn
\teachblock{学生问题诊断}{
  \errortag{盲目代数化简}
  学生拿到此题后，最容易陷入将 $(1-t)a+(1+2t)b+3c=0$ 代入 $\Delta \le 0$ 中进行繁琐的二次方程消元，最后由于多参混合、符号混乱而无法进行下去。
  
  \errortag{忽视切点分支}
  解方程 $4k^2+5k+1=0$ 得到两个斜率值后，很多学生会直接代入较小的 $k=-1$ 计算。\\
  这会导致漏掉对“第一象限切点”这一硬性几何边界的校验，从而得出错解。
}
\teachblock{课堂处理与追问}{
  \textbf{【课堂引入提问口径】}\\
  “我们以前做线性规划，可行域边界都是直线。但这里参数满足：
  $$ac \ge b^2/4$$
  如果把 $a,c$ 看成 $x,y$，这个边界是双曲线！而目标 $t$ 分离后是一个分式，能让我们联想到两点连线的斜率。”
  
  \textbf{【核心追问清单】}\\
  \begin{enumerate}[label=\arabic*., leftmargin=1.5em]
    \item 联立消元得到切率 $k=-1$ 和 $-1/4$。为什么舍去 $k=-1$？它在几何上对应什么？
          （答：它对应定点 $P$ 向双曲线的第三象限分支作的切线，此时切点横坐标 $a<0$，不符合 $a>0$ 的题意。）
    \item 双曲线的开口和顶点由谁决定？（答：由常数 $b^2/4$ 决定，且定点 $P$ 的位置也是由 $b$ 决定的，因此图形具有自相似性。）
  \end{enumerate}
}
\end{paracol}

\newpage

\teacherproblem{二}{偶函数唯一零点与投影点轨迹}

\begin{exbox}{【例题 2】}
若函数 $f(x) = x^2 - a|x| + \frac{\cos x}{x^2+1} + a$ 有且只有一个零点。点 $P(3a, 1)$ 在动直线 $m(x-1) + n(y-1) = 0$ 上的投影为点 $M$（其中直线参数满足 $(m,n) \ne (0,0)$）。若定点 $N(3, 3)$，求 $|MN|$ 的最小值。
\end{exbox}

\par\addvspace{0.4em}{\centering
\begin{tikzpicture}[scale=0.8, >=stealth]
  % 网格
  \draw[color=gridcolor, thin, step=1cm] (-6,-2.5) grid (5,5);
  % 坐标轴
  \draw[->, thick] (-6.5,0) -- (5.5,0) node[right] {$x$};
  \draw[->, thick] (0,-2.8) -- (0,5.5) node[above] {$y$};
  \node[below left] at (0,0) {$O$};
  
  % 圆轨迹 (x+1)^2 + (y-1)^2 = 4
  \coordinate (Center) at (-1,1);
  \draw[thick, blue] (Center) circle (2cm);
  \fill[blue] (Center) circle (1.5pt) node[below] {$O_0(-1,1)$};
  
  % 定点 P(-3, 1) 和 A(1, 1)
  \coordinate (P) at (-3,1);
  \coordinate (A) at (1,1);
  \fill[red] (P) circle (2pt) node[left] {$P(-3,1)$};
  \fill[red] (A) circle (2pt) node[right] {$A(1,1)$};
  \draw[gray, thin, dashed] (P) -- (A);
  
  % 动点 M
  \coordinate (M) at (0, 2.732);
  \fill[black] (M) circle (1.8pt) node[above left] {$M$};
  
  % 动直线 AM
  \draw[thick, orange, domain=-1.5:2.5] plot (\x, {1.732*(\x-1) + 1});
  \node[orange, right] at (2.2, 3.0) {动直线 $l$};
  
  % 垂直线 PM
  \draw[thick, dashed, purple] (P) -- (M);
  
  % 直角标记 at M
  \draw[thin] ($(M)!0.2!(P)$) -- ($($(M)!0.2!(P)$) + ($(M)!0.2!(A)$) - (M)$) -- ($(M)!0.2!(A)$);
  
  % 定点 N(3, 3)
  \coordinate (N) at (3,3);
  \fill[red] (N) circle (2pt) node[above right] {$N(3,3)$};
  
  % 连线 ON
  \draw[thick, green!60!black] (Center) -- (N);
  
  % 轨迹最接近点 M0
  \coordinate (M0) at (0.789, 1.894);
  \fill[red] (M0) circle (1.8pt) node[above left] {$M_0$};
  
  \node[below] at (-1, -1.5) {投影点 $M$ 轨迹是以 $PA$ 为直径的圆};
\end{tikzpicture}\par}\addvspace{0.4em}

\begin{paracol}{2}
\teachblock{标准解答与讲解主线}{
  \textbf{第一步：由偶函数唯一零点确定参数 $a$}\\
  我们首先证明函数 $f(x)$ 的奇偶性。由于自变量 $x$ 仅以平方、绝对值和余弦的形式出现：
  \begin{itemize}
    \item 对于二次项：$(-x)^2 = x^2$；
    \item 对于绝对值项：$|-x| = |x|$；
    \item 对于余弦分式项：$\frac{\cos(-x)}{(-x)^2+1} = \frac{\cos x}{x^2+1}$。
  \end{itemize}
  因此对定义域内的任意实数 $x$ 均有：
  $$f(-x) = (-x)^2 - a|-x| + \frac{\cos(-x)}{(-x)^2+1} + a = f(x)$$
  即该函数为偶函数。\\
  偶函数的图像关于 $y$ 轴对称。若非零实数 $x_0 \ne 0$ 是函数的一个零点（即 $f(x_0) = 0$），则其对称点 $-x_0$ 也必是零点。这意味着偶函数的非零零点必然是成对出现的（一正一负，至少有两个零点）。\\
  题目给定该函数有且仅有一个零点，因此唯一的零点只能是其自身的对称点，即必为 $x = 0$。\\
  将唯一零点 $x = 0$ 代入函数中，必有 $f(0) = 0$：
  $$\begin{aligned}
    f(0) &= 0^2 - a|0| + \frac{\cos 0}{0^2+1} + a = 0 \\
    &\implies 0 - 0 + \frac{1}{1} + a = 0 \\
    &\implies 1 + a = 0 \\
    &\implies a = -1
  \end{aligned}$$
  \textbf{第二步：代入检验唯一零点的充分性}\\
  我们求得 $a = -1$。为了保证逻辑的严密性，必须验证当 $a = -1$ 时，函数是否确实只有 $x=0$ 这一个零点（即当 $x \ne 0$ 时，函数值永远不为 $0$）。\\
  将 $a = -1$ 代入函数解析式：
  $$f(x) = x^2 + |x| + \frac{\cos x}{x^2+1} - 1$$
  我们利用常用不等式放缩关系：对于任意实数 $x$ 均有 $\cos x \ge 1 - x^2/2$。\\
  代入估计 $f(x)$ 的下界：
  $$f(x) \ge x^2 + |x| - 1 + \frac{1 - x^2/2}{x^2+1}$$
  设 $u = |x| > 0$（因为考虑 $x \ne 0$ 的非零情况），则 $x^2 = u^2$。代入进行通分整理：
  $$f(x) \ge u^2 + u - 1 + \frac{1 - u^2/2}{u^2+1} = \frac{(u^2+u-1)(u^2+1) + (1-u^2/2)}{u^2+1}$$
  将分子展开合并同类项：
  $$(u^4 + u^3 - u^2 + u^2 + u - 1) + 1 - u^2/2 = u^4 + u^3 + u - \frac{u^2}{2} = u\left(u^3 + u^2 - \frac{u}{2} + 1\right)$$
  我们将括号内的项进行配方：
  $$u^3 + u^2 - \frac{u}{2} + 1 = u^3 + \left(u - \frac{1}{4}\right)^2 - \frac{1}{16} + 1 = u^3 + \left(u - \frac{1}{4}\right)^2 + \frac{15}{16}$$
  因为 $u > 0 \implies u^3 > 0$，且二次项非负，常数项大于零，所以括号内的表达式必大于零。\\
  既然分子大于 $0$ 且分母 $u^2+1 > 0$，所以对于任意 $x \ne 0$ 均有 $f(x) > 0$。\\
  故 $a = -1$ 时，函数确实有且仅有一个零点 $x = 0$。充分性验证成立。\\[0.3em]
  \textbf{第三步：确定定点 $P$ 与直线定点 $A$ 的坐标}\\
  由 $a = -1$ 得到定点 $P(3a, 1) = P(-3, 1)$。\\
  动直线方程为 $m(x-1) + n(y-1) = 0$。我们令 $x-1=0, y-1=0$，解得 $x=1, y=1$。\\
  这说明不论直线参数 $m,n$ 取何值，只要代入 $x=1, y=1$，等式左边必为 $m(0) + n(0) = 0$，等式恒成立。\\
  因此动直线恒过定点 $A(1,1)$。\\[0.3em]
  \textbf{第四步：确定投影点 $M$ 的轨迹圆}\\
  点 $M$ 是点 $P$ 在经过定点 $A$ 的动直线上的投影，由投影定义可知，线段 $PM$ 与该动直线垂直，即：
  $$PM \perp AM \implies \angle PMA = 90^\circ$$
  由平面几何定理（直角三角形斜边中线等于斜边一半），若设线段 $PA$ 的中点为 $O_0$，则有 $O_0M = \frac{1}{2}PA = \text{常数}$。\\
  这说明动点 $M$ 到定点 $O_0$ 的距离是恒定的。因此点 $M$ 的运动轨迹是以线段 $PA$ 为直径的圆。\\
  圆心 $O_0$ 即为线段 $PA$ 的中点坐标：
  $$O_0 = \left(\frac{-3+1}{2},\ \frac{1+1}{2}\right) = (-1, 1)$$
  半径为 $r = \frac{1}{2}|PA| = \frac{1 - (-3)}{2} = 2$。\\
  因为动直线可绕 $A$ 旋转任意方向，圆上的每一个点都可由相应的动直线投影得到，因此轨迹是完整的圆。\\[0.3em]
  \textbf{第五步：求 $|MN|$ 的最小值}\\
  定点为 $N(3,3)$。计算定点 $N$ 到圆心 $O_0(-1,1)$ 的距离：
  $$O_0N = \sqrt{[3 - (-1)]^2 + (3 - 1)^2} = \sqrt{4^2 + 2^2} = \sqrt{20} = 2\sqrt{5}$$
  因为 $O_0N = 2\sqrt{5} \approx 4.47 > 2 = r$，所以点 $N$ 位于圆的外侧。\\
  由几何直观可知，圆外定点到圆上动点的最小距离，即为圆心到该定点连线段与圆的交点 $M_0$ 对应的长度。\\
  所以，最小距离为：
  $$|MN|_{\min} = O_0N - r = 2\sqrt{5} - 2$$
}
\switchcolumn
\teachblock{学生问题诊断}{
  \errortag{混淆奇偶性质}
  部分学生无法快速识别该分式项的奇偶性，或者在求得 $a = -1$ 后，遗漏了“充分性检验”这一关键步骤。在高考阅卷中，不检验唯一性通常会被扣分。
  
  \errortag{硬联立求垂足}
  对于投影点，学生很容易使用点到直线的垂足公式，将 $M$ 的坐标写成关于 $m,n$ 的参数方程，然后代入距离公式。这种纯代数算术不仅运算量极大，而且极难消去参数 $m,n$ 得到最值。
}
\teachblock{课堂处理与追问}{
  \textbf{【核心讲解引导】}\\
  “求投影，就是找垂直。看到垂直，大家脑海里要蹦出一个什么图形？对，是圆！如果一个角永远是直角，它的顶点轨迹就是一个以斜边为直径的圆。这道题我们不需要知道垂足的坐标，只需要知道它是圆上的一个动点！”
  
  \textbf{【核心追问清单】}\\
  \begin{enumerate}[label=\arabic*., leftmargin=1.5em]
    \item 偶函数若有且仅有三个零点，它们的分布有什么规律？（答：其中一个必为 $0$，另外两个互为相反数。）
    \item 如果直线方程限制为 $y-1=k(x-1)$，投影点 $M$ 的轨迹会有变化吗？
          （答：有。此时若直线斜率存在，则无法表示垂直于 $x$ 轴的直线 $x=1$。因此，圆轨迹上会缺失极限位置点。本题直线族满足 $(m,n)\ne(0,0)$，因此轨迹圆是完整的。）
  \end{enumerate}
}
\end{paracol}

\newpage

\teacherproblem{三}{根式函数的三角形面积模型}

\begin{exbox}{【例题 3】}
已知 $a$ 为常数，函数 $f(x) = \frac{x\left(\sqrt{a-x^2} + \sqrt{1-x^2}\right)}{a-1}$ 的最大值为 $1$，求 $a$ 的所有取值。
\end{exbox}

\par\addvspace{0.45em}{\centering
\begin{minipage}[t]{0.48\textwidth}
\begin{tikzpicture}[scale=1.5, >=stealth]
  % 顶点
  \coordinate (H) at (0,0);
  \coordinate (A) at (0,1.2);
  \coordinate (B) at (-1.6,0);
  \coordinate (C) at (1.0,0);
  
  % 画三角形
  \draw[thick] (A) -- (B) -- (C) -- cycle;
  % 画高 AH
  \draw[thick, dashed, blue] (A) -- (H);
  
  % 标点
  \fill (A) circle (1.2pt) node[above] {$A$};
  \fill (B) circle (1.2pt) node[below left] {$B$};
  \fill (C) circle (1.2pt) node[below right] {$C$};
  \fill (H) circle (1.2pt) node[below] {$H$};
  
  % 标边长
  \node[above left] at ($(A)!0.5!(B)$) {$AB = \sqrt{a}$};
  \node[above right] at ($(A)!0.5!(C)$) {$AC = 1$};
  \node[left, blue] at ($(A)!0.5!(H)$) {$x$};
  \node[below] at ($(B)!0.5!(H)$) {\small $\sqrt{a-x^2}$};
  \node[below] at ($(H)!0.5!(C)$) {\small $\sqrt{1-x^2}$};
  
  % 直角标记 at H
  \draw[thin] (-0.12,0) -- (-0.12,0.12) -- (0,0.12);
  
  \node[below=0.5cm, align=center] at (0,0) {图 1：一般状态下的三角形拼接\\ 底边 $BC = \sqrt{a-x^2} + \sqrt{1-x^2}$};
\end{tikzpicture}
\end{minipage}
\hfill
\begin{minipage}[t]{0.48\textwidth}
\begin{tikzpicture}[scale=1.5, >=stealth]
  % 顶点
  \coordinate (A2) at (0,1.2);
  \coordinate (B2) at (-1.414,1.2); 
  \coordinate (C2) at (0,0.2);
  
  % 画直角三角形
  \draw[thick] (A2) -- (B2) -- (C2) -- cycle;
  
  % 标点
  \fill (A2) circle (1.2pt) node[above right] {$A$};
  \fill (B2) circle (1.2pt) node[left] {$B$};
  \fill (C2) circle (1.2pt) node[below] {$C$};
  
  % 标边长
  \node[above] at ($(A2)!0.5!(B2)$) {$AB = \sqrt{a}$};
  \node[right] at ($(A2)!0.5!(C2)$) {$1$};
  
  % 直角标记 at A
  \draw[thin] (-0.12,1.2) -- (-0.12,1.08) -- (0,1.08);
  
  \node[below=1.5cm, align=center] at (-0.7,0.2) {图 2：最大面积临界状态\\ $\angle BAC = 90^\circ$ 时，$S_{\max} = \frac{1}{2}\sqrt{a}$};
\end{tikzpicture}
\end{minipage}\par}\addvspace{0.45em}

\begin{paracol}{2}
\teachblock{标准解答与讲解主线}{
  \textbf{第一步：定义域与参数范围分析}\\
  为使根式函数在实数范围内有定义，各根式内部表达式必须非负：
  $$a - x^2 \ge 0 \implies x^2 \le a$$
  $$1 - x^2 \ge 0 \implies x^2 \le 1$$
  为了使定义域不为空集（即存在实数 $x$ 使得上述不等式同时成立），参数 $a$ 必须满足：
  $$a \ge 0$$
  同时，因为原解析式中分母为 $a - 1$，分母不能为零，因此 $a \ne 1$。\\
  如果 $a = 0$，则上述不等式限制为 $x^2 \le 0$ 且 $x^2 \le 1$，解得定义域只有唯一一个点 $x=0$。此时自变量只能取 $0$，代入求得唯一函数值 $f(0) = 0$，最大值不可能为 $1$，舍去。\\
  综上，参数 $a$ 的取值范围必须满足：$a > 0$ 且 $a \ne 1$。\\[0.3em]
  \textbf{第二步：利用奇偶性简化分类讨论}\\
  我们观察函数分子部分为 $g(x) = x\left(\sqrt{a-x^2} + \sqrt{1-x^2}\right)$。对任意定义域内的 $x$ 均有：
  $$g(-x) = -x\left(\sqrt{a-(-x)^2} + \sqrt{1-(-x)^2}\right) = -g(x)$$
  故分子为奇函数。而分母 $a-1$ 为非零常数，所以 $f(x)$ 整体是定义域上的奇函数。\\
  根据奇函数的对称性，其正半轴与负半轴的函数值互为相反数。如果它在正半轴的最大值为 $M$，则在负半轴的最小值为 $-M$。\\
  因此，我们只需要研究正半轴 $x \ge 0$ 时分子的最大值 $g_{\max}$，然后再根据分母 $a-1$ 的正负号来判断最大值具体是在正半轴还是负半轴取得。\\[0.3em]
  \textbf{第三步：构造直角三角形几何拼接模型}\\
  对于 $x \ge 0$，我们构造如图 1 所示的三角形 $ABC$。自顶点 $A$ 向底边 $BC$ 引高线 $AH$，使高为 $AH = x$。\\
  使垂足 $H$ 将底边 $BC$ 分为两部分，其长度分别为：
  $$BH = \sqrt{a-x^2},\qquad HC = \sqrt{1-x^2}$$
  根据直角三角形的勾股定理，两斜边的长度分别为：
  $$AB = \sqrt{AH^2 + BH^2} = \sqrt{x^2 + (a - x^2)} = \sqrt{a}$$
  $$AC = \sqrt{AH^2 + HC^2} = \sqrt{x^2 + (1 - x^2)} = 1$$
  由于 $AB$ 和 $AC$ 的长度仅由常数 $a$ 决定，因此不论 $x$ 如何变化，三角形 $ABC$ 的侧边长是固定不变的。\\
  此时，底边长度为两段之和：$BC = BH + HC = \sqrt{a-x^2} + \sqrt{1-x^2}$。\\
  则三角形 $ABC$ 的面积为：
  $$S_{\triangle ABC} = \frac{1}{2} BC \cdot AH = \frac{1}{2} x\left(\sqrt{a-x^2} + \sqrt{1-x^2}\right)$$
  对比分子解析式 $g(x)$，可以发现分子恰好是三角形面积的两倍：
  $$g(x) = 2S_{\triangle ABC}$$
  \textbf{第四步：由正弦面积公式确定分子最大值}\\
  设两条固定侧边 $AB$ 与 $AC$ 的夹角为 $\theta$（如图 2），则该三角形的面积为：
  $$S_{\triangle ABC} = \frac{1}{2} AB \cdot AC \cdot \sin\theta = \frac{1}{2} \cdot \sqrt{a} \cdot 1 \cdot \sin\theta = \frac{\sqrt{a}}{2}\sin\theta$$
  因为对于任意三角形的夹角，均有正弦值 $\sin\theta \le 1$。\\
  所以，当且仅当 $\theta = 90^\circ$（即两侧边 $AB$ 与 $AC$ 互相垂直）时，三角形面积取得最大值：
  $$S_{\max} = \frac{\sqrt{a}}{2}$$
  此时分子取得的最大值为：
  $$g_{\max} = 2S_{\max} = \sqrt{a}$$
  \textbf{第五步：等号成立条件的可行性校验}\\
  我们需要验证 $\theta = 90^\circ$ 的直角状态是否能在自变量定义域内取得。\\
  当 $\angle BAC = 90^\circ$ 时，由直角三角形的射影定理可知，斜边上的高平方等于两投影段乘积：
  $$AH^2 = BH \cdot HC \implies x^2 = \sqrt{a-x^2}\sqrt{1-x^2}$$
  两边平方整理：
  $$\begin{aligned}
    x^4 &= (a-x^2)(1-x^2) \\
    &\implies x^4 = a - (a+1)x^2 + x^4 \\
    &\implies (a+1)x^2 = a \\
    &\implies x^2 = \frac{a}{a+1}
  \end{aligned}$$
  因为已知 $a > 0$，所以分母 $a+1 > a > 0$，从而有：
  $$0 < x^2 = \frac{a}{a+1} < 1 \quad \text{且} \quad x^2 = \frac{a}{a+1} < a$$
  这说明求出的 $x^2$ 同时满足 $x^2 \le 1$ 且 $x^2 \le a$，即等号成立时的自变量值完全落在定义域内。\\
  因此分子最大值 $g_{\max} = \sqrt{a}$ 确实可以取得。\\[0.3em]
  \textbf{第六步：根据分母符号分类讨论解出参数 $a$}
  \begin{itemize}
    \item \textbf{情形一：当 $a > 1$ 时}，分母 $a-1 > 0$。\\
          此时，函数最大值在分子取最大值时取得（此时 $x > 0$）：
          $$f_{\max} = \frac{g_{\max}}{a-1} = \frac{\sqrt{a}}{a-1}$$
          已知函数最大值为 $1$，得方程：
          $$\frac{\sqrt{a}}{a-1} = 1 \implies \sqrt{a} = a - 1$$
          设 $u = \sqrt{a} > 1$（因为 $a > 1$），则方程为 $u = u^2 - 1 \implies u^2 - u - 1 = 0$。\\
          解得 $u = \frac{1 \pm \sqrt{5}}{2}$。舍去小于零的负根，保留正根 $u = \frac{1+\sqrt{5}}{2}$。\\
          从而求得：
          $$a = u^2 = \left(\frac{1+\sqrt{5}}{2}\right)^2 = \frac{3+\sqrt{5}}{2}$$
          由于 $\frac{3+\sqrt{5}}{2} \approx 2.618 > 1$，符合情形一的限制。
    \item \textbf{情形二：当 $0 < a < 1$ 时}，分母 $a-1 < 0$。\\
          由奇函数对称性，由于 $x \ge 0$ 时函数值均非正，最大值必在 $x < 0$ 负半轴处取得，其大小为：
          $$f_{\max} = \frac{g_{\max}}{1-a} = \frac{\sqrt{a}}{1-a}$$
          已知函数最大值为 $1$，得方程：
          $$\frac{\sqrt{a}}{1-a} = 1 \implies \sqrt{a} = 1 - a$$
          设 $u = \sqrt{a}$ 且 $0 < u < 1$（因为 $0 < a < 1$），则方程为 $u = 1 - u^2 \implies u^2 + u - 1 = 0$。\\
          解得 $u = \frac{-1 \pm \sqrt{5}}{2}$。舍去不符合区间的根，保留 $u = \frac{\sqrt{5}-1}{2}$。\\
          从而求得：
          $$a = u^2 = \left(\frac{\sqrt{5}-1}{2}\right)^2 = \frac{3-\sqrt{5}}{2}$$
          由于 $\frac{3-\sqrt{5}}{2} \approx 0.382 \in (0,1)$，符合情形二的限制。
  \end{itemize}
  综上所述，常数 $a$ 的所有取值为：$a = \frac{3-\sqrt{5}}{2}$ 或 $a = \frac{3+\sqrt{5}}{2}$。
}
\switchcolumn
\teachblock{学生问题诊断}{
  \errortag{导数运算爆炸}
  如果直接对该函数求导，导数解析式中将包含四个根式乘积及分子分母的复合，极易产生庞大的运算量，绝大多数学生在考场上根本无法算出正确结果。
  
  \errortag{漏掉 $0<a<1$ 分支}
  学生极易忽视分母 $a-1$ 的符号。如果直接解方程 $\sqrt{a}/(a-1)=1$，会得出唯一解 $a = (3+\sqrt{5})/2$，从而彻底遗漏了 $0<a<1$ 时的另一个解。这主要是因为缺乏对奇函数在负半轴取得最大值的对称性思考。
}
\teachblock{课堂处理与追问}{
  \textbf{【几何模型启发】}\\
  “如果把 $\sqrt{a-x^2}$ 和 $\sqrt{1-x^2}$ 看成两条直角边，高为 $x$，我们可以拼出什么？对，两个直角三角形拼成一个大三角形。拼出来的三角形，有两条边是固定的：$\sqrt{a}$ 和 $1$。想一想，两条边固定时，三角形什么时候面积最大？当然是它们垂直夹角为 90 度的时候！”
  
  \textbf{【核心追问清单】}\\
  \begin{enumerate}[label=\arabic*., leftmargin=1.5em]
    \item 为什么我们可以直接说最大值是 $\sqrt{a}/|a-1|$？
          （答：因为函数是奇函数，如果分母为正，最大值在正半轴取得；如果分母为负，最大值在负半轴由奇函数性质变为等价的绝对值分母形式。）
    \item 在拼图的过程中，我们假设了 $BH$ 和 $HC$ 在高的两侧（即底边长为两段之和）。如果它们在高的同侧，对应的代数式是什么？\\
          （答：对应的是两根式之差：
          $$|\sqrt{a-x^2} - \sqrt{1-x^2}|$$
          这在后续变式中会提到。）
  \end{enumerate}
}
\end{paracol}

\newpage

% ==================== 第三部分 ====================
\section{应试方法联系与深度思考}

\subsection{1. 方法的横向联系}

\begin{itemize}
  \item \textbf{1.1 “分式与斜率”的横向迁移}：\\
        形如 $(y-b)/(x-a)$ 的代数式都可以解释为平面动点 $(x,y)$ 与定点 $(a,b)$ 的连线斜率。这一思想不仅可用于本课第一题，还可广泛迁移到：
        \begin{enumerate}[label=(\arabic*)]
          \item 圆锥曲线中弦的斜率范围问题；
          \item 导数中割线斜率的临界制造与中值定理几何解释；
          \item 形如 $(Ax+By+C)/(Dx+Ey+F)$ 的双线性分式最值问题（转化为定点到可行域的切线）。
        \end{enumerate}
  \item \textbf{1.2 “恒成立与可行域”的非线性映射}：\\
        判别式约束 $\Delta \le 0$ 不仅是一个代数不等式，当参数多于两个时，它将确定一个平面非线性区域。将参数看作坐标，可行域的边界（如双曲线、圆或抛物线）就是代数相切的几何临界线。这对于理解高等数学中的包络线、约束最值有极大帮助。
  \item \textbf{1.3 “投影与直径圆”的几何本质}：\\
        “投影 $=$ 垂直 $=$ 直角 $=$ 直径圆”。在解析几何中，看到“点 $M$ 在动直线上，且满足某种垂直关系”，或“向量数量积为 0”，要立刻摆脱代数联立，用定线段为直径的圆轨迹直接写出动点轨迹方程。
  \item \textbf{1.4 “根式与直角三角形面积”的无导数化简}：\\
        双根式复合代数式 $\sqrt{R^2 - x^2}$ 具有天然的勾股几何背景。通过将自变量看作高，将根式看作底边分段，将复杂的复合函数最值转化为面积最值，从而避开求导，这是化归与转化思想的巅峰体现。
\end{itemize}

\subsection{2. 解题策略取舍：通法与巧法}

\begin{center}
\renewcommand{\arraystretch}{1.3}
\begin{tabularx}{0.98\linewidth}{l p{4.2cm} p{4.2cm} X}
\toprule
\textbf{题型} & \textbf{常规代数通法} & \textbf{几何直观巧法} & \textbf{方法对比与适用边界} \\
\midrule
\textbf{第一题} & 研究目标式关于自变量 $b$ 的单调性，代入最值后换元配方。 & 转化为双曲线可行域外侧与定点连线的斜率最值。 & \textbf{对比}：代数法单调性讨论复杂，极易算错；几何法临界相切极快。\\
& & & \textbf{边界}：分式必须能化为斜率结构，且参数可行域边界较规则。 \\
\midrule
\textbf{第二题} & 设直线斜率，求出垂足坐标表达式，代入两点距离公式求最值。 & 识别直线恒过定点，利用直角投影锁定圆轨迹，用圆心距减半径。 & \textbf{对比}：代数法求垂足坐标极其繁琐；几何法无需知道垂足即可破局。\\
& & & \textbf{边界}：直线族必须恒过定点，否则垂足轨迹不为圆。 \\
\midrule
\textbf{第三题} & 强行求导，通分后讨论分子零点及参数分类。 & 构造直角三角形，将分子转化为面积，由夹角为直角时面积最大求解。 & \textbf{对比}：求导运算量属于灾难级；几何法 3 行得出分子最大值。\\
& & & \textbf{边界}：根式内部必须呈平方差勾股结构，且必须校验等号点在定义域内。 \\
\bottomrule
\end{tabularx}
\end{center}

\subsection{3. 命题视角的变式与拓展}

\begin{itemize}
  \item \textbf{3.1 第一题的边界变式（双曲线改圆）}：\\
        若参数约束条件改为 $a^2 + c^2 \le R^2$，目标式仍为 $t = 1 + 3(b+c)/(a-2b)$。\\
        此时，动点 $Q(a,c)$ 的可行域变为圆盘内部。定点 $P(2b,-b)$ 到圆盘的斜率最值同样在切线处取得，可通过点到直线的距离公式 $d = |A x_0 + B y_0 + C|/\sqrt{A^2+B^2} = R$ 快速列方程求解，避开复杂的消元判别式。
  \item \textbf{3.2 第二题的夹角变式（直角改定角）}：\\
        若将“投影点（垂直）”改为“$\angle PMA = 60^\circ$”，则点 $M$ 的轨迹将不再是以 $PA$ 为直径的圆，而是以 $PA$ 为弦、所对圆周角为 $60^\circ$ 的两段圆弧（即著名的“阿波罗尼斯圆弧”）。轨迹仍然是圆的一部分，可用圆的几何性质求解。
  \item \textbf{3.3 第三题的运算变式（加号改减号）}：\\
        若分子函数改为 $g(x) = x\left(\sqrt{a-x^2} - \sqrt{1-x^2}\right)$。\\
        在几何上，这对应高 $AH=x$ 落在底边 $BC$ 的延长线上（即 $H$ 在 $BC$ 外部），底边长度为两段直角边之差 $|BH - HC|$。此时，原三角形拼接面积模型需要修正为大三角形与小三角形的面积差，最值不再是简单的 $\sqrt{a}$，需重新建立函数关系。
  \item \textbf{3.4 第三题的输出变式（最大值改值域）}：\\
        若题目要求函数 $f(x)$ 的值域，由于函数在对称定义域内是奇函数，若求得其最大值为 $M = \sqrt{a}/|a-1|$，则其最小值为 $-M$。\\
        在 $a > 0$ 且 $a \ne 1$ 时，其值域可直接写为：
        $$\left[-\frac{\sqrt{a}}{|a-1|},\ \frac{\sqrt{a}}{|a-1|}\right]$$
\end{itemize}

\subsection{4. 实战做题启发}

\begin{keybox}{实战启发总结}
\begin{enumerate}[label=\arabic*., leftmargin=2em, itemsep=0.2em]
  \item \textbf{代数式的几何雷达}：\\
        看到一次分式优先想斜率，看到数量积为零或垂足优先想直径圆，看到双根式平方差优先想直角三角形边长。
  \item \textbf{先找定点，以静制动}：\\
        处理含有多参数的直线、圆锥曲线时，第一动作是寻找“恒过的定点”，这能瞬间明确图形旋转或伸缩的几何核心。
  \item \textbf{必须进行边界与等号点校验}：\\
        几何巧法求出最值后，必须返回代数定义域，检验临界等号点（如相切时的切点横坐标、直角三角形高 $x^2 = \frac{a}{a+1}$）是否在定义域内，以及切点是否在合法的曲线上（排除虚假分支）。
  \item \textbf{奇偶性与对称性是分类讨论的减震器}：\\
        在求复合代数式最值前，先判断函数的奇偶性。利用对称性将定义域缩减为正半轴，求出最值后再通过绝对值或正负号还原，可极大简化分类讨论。
\end{enumerate}
\end{keybox}

\end{document}
